\label{der:sylvester-full}
复数方阵可上三角化：取一个特征向量补成基，首列只在首位有特征值，右下块再归纳，得到$A=PTP\inv,B=QSQ\inv$，其中$T,S$上三角。设$Y=P\inv XQ,D=P\inv CQ$，原方程等价于$TY-YS=D$。

逐元素写出第$(i,j)$式：
\[(t_{ii}-s_{jj})y_{ij}=d_{ij}-\sum_{k=i+1}^m t_{ik}y_{kj}+\sum_{k=1}^{j-1}y_{ik}s_{kj}.\]
按$j=1,2,\ldots,n$逐列，并在每列按$i=m,m-1,\ldots,1$向上求。右边第一和只用当前列更下面已求出的元素，第二和只用以前各列已求出的元素。谱不交保证每个$t_{ii}-s_{jj}\ne0$，因此此顺序唯一算出全部$mn$个元素。将这些元素拼成$Y$，每条方程按构造成立，再取$X=PYQ\inv$得到存在解；任何解都必须依此顺序逐项相同，故唯一。

若$A,B$已对角化，两和均为零，简化为$y_{ij}=\frac{d_{ij}}{\lambda_i-\mu_j}$。实系数情形虽然允许在复数域求解，但共轭后的解仍满足同一方程，由唯一性知解等于自身共轭，所以最后解为实矩阵。
